% arrow_comparison.tex — Self-contained section for monograph
% \input this into the main document; requires references.bib entries for
% arrow1951social, nipkow2009social, okasha2011theory, gibbard1973manipulation,
% satterthwaite1975strategy.

% Section heading is in the parent file (universal_impossibility_monograph.tex)

Two impossibility theorems constrain aggregation under conflicting
desiderata: Arrow's impossibility theorem in social choice
\citep{arrow1951social} and the universal explanation impossibility
developed in this work. Both assert that a small set of individually
reasonable properties cannot be jointly satisfied. This section maps
the axioms of each theorem onto the other, identifies the precise point
where the correspondence breaks, and characterizes the complementary
regimes in which the two results operate.

%% ===================================================================
\subsection{Formal setup}
\label{sec:arrow-setup}

\paragraph{Arrow's theorem.}
Let $A$ be a finite set of alternatives with $|A| \geq 3$, and let
$N = \{1, \ldots, n\}$ with $n \geq 2$ be a set of voters. Write
$\mathcal{L}(A)$ for the set of strict linear orders on~$A$ (complete,
transitive, antisymmetric, irreflexive). A \emph{preference profile}
is a tuple $\pi = (\pi_1, \ldots, \pi_n) \in \mathcal{L}(A)^N$.
A \emph{social welfare function} is a map
$F : \mathcal{L}(A)^N \to \mathcal{L}(A)$. Arrow's theorem states
that no~$F$ satisfies all four of the following axioms:
\begin{enumerate}[label=\textbf{A\arabic*.},leftmargin=2.5em]
  \item \textbf{Unrestricted Domain (UD).}\; $F$ is defined on all
    of~$\mathcal{L}(A)^N$.
  \item \textbf{Weak Pareto (WP).}\; If $a \succ_{\pi_i} b$ for every
    voter~$i$, then $a \succ_{F(\pi)} b$.
  \item \textbf{Independence of Irrelevant Alternatives (IIA).}\;
    For all $a, b \in A$: if profiles $\pi$ and $\pi'$ agree on every
    voter's ranking of~$a$ versus~$b$, then $F(\pi)$ and $F(\pi')$
    agree on~$a$ versus~$b$.
  \item \textbf{Non-Dictatorship (ND).}\; There is no voter~$i$ such
    that $F(\pi) = \pi_i$ for all~$\pi$.
\end{enumerate}

\paragraph{The explanation impossibility.}
An \textsc{ExplanationSystem}
$(\Theta, H, Y, \mathrm{obs}, \mathrm{exp}, \mathrm{incomp})$
consists of a configuration space~$\Theta$, an explanation
space~$H$, an observable space~$Y$, an observation map
$\mathrm{obs} : \Theta \to Y$, a native explanation map
$\mathrm{exp} : \Theta \to H$, and an irreflexive incompatibility
relation on~$H$. We say the system has the \emph{Rashomon property}
if there exist $\theta_1, \theta_2 \in \Theta$ with
$\mathrm{obs}(\theta_1) = \mathrm{obs}(\theta_2)$ and
$\mathrm{incomp}(\mathrm{exp}(\theta_1),
\mathrm{exp}(\theta_2))$. For any such system, no explanation
$E : \Theta \to H$ is simultaneously:
\begin{enumerate}[label=\textbf{E\arabic*.},leftmargin=2.5em]
  \item \textbf{Faithful.}\;
    $\neg\,\mathrm{incomp}(E(\theta),\,\mathrm{explain}(\theta))$
    for all~$\theta$.
  \item \textbf{Stable.}\; $\mathrm{observe}(\theta_1) =
    \mathrm{observe}(\theta_2) \Rightarrow E(\theta_1) = E(\theta_2)$.
  \item \textbf{Decisive.}\;
    $\mathrm{incomp}(\mathrm{explain}(\theta), h)
    \Rightarrow \mathrm{incomp}(E(\theta), h)$ for all~$\theta, h$.
\end{enumerate}
This is \texttt{explanation\_impossibility} in the Lean formalization,
proved from the Rashomon property alone with zero auxiliary axioms.

%% ===================================================================
\subsection{The axiom correspondence}
\label{sec:arrow-axiom-map}

We now map each Arrow axiom to the explanation framework. The mapping
uses the \emph{query-relative} formulation (Lean:
\texttt{QueryRelative.lean}), in which incompatibility is parameterized
by a query space~$Q$, with each query $q \in Q$ inducing its own
incompatibility relation $\mathrm{incomp}_q$ on~$H$.

\begin{table}[ht]
\centering
\renewcommand{\arraystretch}{1.3}
\begin{tabular}{@{}lllc@{}}
\toprule
\textbf{Arrow axiom} & \textbf{Explanation property} & \textbf{Quality} \\
\midrule
IIA & Query-relative stability & Exact \\
Weak Pareto & Faithfulness (restricted) & Partial \\
Non-Dictatorship & (no analog) & None \\
Unrestricted Domain & (implicit totality) & Structural \\
\bottomrule
\end{tabular}
\caption{Axiom correspondence between Arrow's theorem and the explanation
impossibility. Only IIA admits an exact translation.}
\label{tab:arrow-map}
\end{table}

\paragraph{IIA \boldmath$\longleftrightarrow$ query-relative stability
  (exact).}
Fix a pair $q = (a,b) \in A \times A$ as a query. Define the
query-relative observation map
$\mathrm{observe}_q(\pi) = (\mathbf{1}[a \succ_{\pi_i} b])_{i \in N}$,
the vector of individual pairwise preferences restricted to~$(a,b)$.
IIA states: if $\mathrm{observe}_q(\pi) = \mathrm{observe}_q(\pi')$,
then the social ranking of~$a$ versus~$b$ agrees under $F(\pi)$ and
$F(\pi')$. This is precisely query-relative stability
(Lean: \texttt{q\_faithful}, \texttt{q\_decisive} in
\texttt{QueryRelative.lean}) applied to the pairwise query~$q$:
$E$'s answer to ``does $a$ rank above~$b$?''\ depends only on the
pairwise observations. The correspondence is exact; no content is
lost or added.

\paragraph{Weak Pareto \boldmath$\longleftrightarrow$ faithfulness
  (partial).}
WP says: if every voter ranks $a \succ b$, then $F$ ranks $a \succ b$.
Faithfulness says: $E(\theta)$ never contradicts
$\mathrm{explain}(\theta)$. The mapping is indirect. If we set
$\mathrm{explain}(\pi) = F(\pi)$, then faithfulness requires that
$E$ not contradict~$F$. If $F$ satisfies WP, faithfulness of~$E$
to~$F$ automatically respects Pareto---but only transitively.
Pareto constrains~$F$ directly (unanimous preferences are respected);
faithfulness constrains~$E$ relative to~$F$. This is a genuine
structural difference: Pareto is a first-order constraint on the
aggregation map, while faithfulness is a second-order constraint on
an \emph{explanation of} the aggregation map.

\paragraph{Non-Dictatorship (no analog).}
ND constrains the internal structure of~$F$: no single voter
determines the output. The explanation framework has no comparable
axiom. Our three properties---faithful, stable, decisive---constrain
an external explanation~$E$ relative to a given system~$S$; they
say nothing about the internal structure of the system's own
maps. In Arrow's framework, $F$ is \emph{both} the system
and the object being evaluated. In our framework, $\mathrm{explain}$
is the system's native output, and $E$ is a separate map that
interprets it. The $E/F$ split has no counterpart in social choice.

\paragraph{Unrestricted Domain (structural).}
Both frameworks implicitly assume totality: $F$ is defined on all
profiles, and $E$ is defined on all configurations. UD plays the same
structural role in each---it ensures the impossibility is not
circumvented by restricting the domain.

%% ===================================================================
\subsection{Why the embedding fails}
\label{sec:arrow-embedding-fails}

The axiom correspondence of Section~\ref{sec:arrow-axiom-map}
suggests a natural question: is Arrow's theorem an instance of the
explanation impossibility? The answer is no, and the reason is precise.

\begin{quote}
\textbf{Key insight.}\; \emph{IIA prevents the Rashomon property from
holding.}
\end{quote}

To apply the explanation impossibility, we need the Rashomon property:
two configurations $\theta_1, \theta_2$ with the same observation
but incompatible explanations. In the social choice mapping with
query $q = (a,b)$, the observation is
$\mathrm{observe}_q(\pi) = (\mathbf{1}[a \succ_{\pi_i} b])_{i \in N}$
and the explanation is the social ranking of~$a$ versus~$b$ under~$F(\pi)$.
Rashomon at query~$q$ would require two profiles $\pi, \pi'$ with
$\mathrm{observe}_q(\pi) = \mathrm{observe}_q(\pi')$ but
$F(\pi)$ and $F(\pi')$ disagreeing on~$(a,b)$.

But this is \emph{exactly what IIA forbids}. If $F$ satisfies IIA,
then $\mathrm{observe}_q(\pi) = \mathrm{observe}_q(\pi')$ implies
that $F(\pi)$ and $F(\pi')$ agree on~$(a,b)$. The social ranking
of any pair is a deterministic function of the individual pairwise
preferences. There is no Rashomon ambiguity for any query: the
observation uniquely determines the explanation.

The two theorems therefore operate in \emph{complementary regimes}:

\begin{center}
\renewcommand{\arraystretch}{1.2}
\begin{tabular}{@{}ll@{}}
\toprule
\textbf{Regime} & \textbf{Impossibility} \\
\midrule
Rashomon present, IIA not assumed &
  No $E$ is faithful + stable + decisive \\
IIA enforced, no pairwise Rashomon &
  Only dictatorships satisfy WP + ND \\
\bottomrule
\end{tabular}
\end{center}

The explanation impossibility says: when the map from observations to
explanations is many-to-one (Rashomon), you cannot explain faithfully,
stably, and decisively. Arrow's theorem says: when the map from
observations to the social ranking is \emph{required} to be
one-to-one per query (IIA), the only aggregation rules that also
respect unanimity are dictatorial. Different failure modes of
aggregation, activated by opposite structural assumptions.

This complementarity is not a deficiency of either theorem. It
reflects a genuine difference in what each result constrains.
The explanation impossibility is about \emph{underspecification}:
multiple models fit the data, and no single explanation can
faithfully represent them all. Arrow's theorem is about
\emph{preference diversity}: multiple voters have conflicting
orderings, and no aggregation rule can satisfy a short list of
fairness criteria. Underspecification and preference diversity are
distinct sources of impossibility, and neither subsumes the other.

%% ===================================================================
\subsection{What the theorems share}
\label{sec:arrow-shared-structure}

Despite the embedding failure, the theorems share substantial
structure:

\begin{enumerate}
  \item \textbf{Both are impossibility results about aggregation.}\;
    Arrow aggregates individual preferences into a social ordering;
    the explanation impossibility aggregates a Rashomon set of models
    into a single explanation.

  \item \textbf{Both have three independent desiderata.}\; Once
    unrestricted domain (or totality of~$E$) is factored out as a
    structural assumption, Arrow has WP, IIA, and ND; the explanation
    impossibility has faithfulness, stability, and decisiveness.

  \item \textbf{Both are tight.}\; Dropping any single axiom makes
    the remaining set satisfiable. For Arrow: dropping ND yields
    dictatorship; dropping IIA yields Borda count; dropping WP yields
    anti-dictatorships. For the explanation impossibility: dropping
    decisiveness yields $G$-invariant projections (e.g., CPDAG, DASH);
    dropping stability yields the native explanation; dropping
    faithfulness yields an arbitrary constant.

  \item \textbf{Both have the same proof shape.}\; Short diagonal
    arguments. Arrow's proof (in the Geanakoplos version) constructs
    a pivotal voter by toggling an alternative between the top and
    bottom of one voter's ranking. The explanation impossibility
    constructs a contradiction by chaining decisiveness, stability,
    and faithfulness through two Rashomon-equivalent configurations.
    Both are four-to-five-step proofs.

  \item \textbf{Both have been formalized in proof assistants.}\;
    \citet{nipkow2009social} formalized Arrow's theorem and
    the Gibbard--Satterthwaite theorem \citep{gibbard1973manipulation,
    satterthwaite1975strategy} in Isabelle/HOL.
    The explanation impossibility is formalized in Lean~4 (this work),
    with zero axiom dependencies.

  \item \textbf{Both involve a map from a product space to a
    codomain, where natural conditions are jointly unsatisfiable.}\;
    Arrow: $F : \mathcal{L}(A)^N \to \mathcal{L}(A)$, conditions
    WP + IIA + ND. Explanation: $E : \Theta \to H$, conditions
    faithful + stable + decisive. The abstract shape is the same:
    a function from a structured input to a structured output,
    subject to three constraints that form an inconsistent triad.
\end{enumerate}

The shared structure is categorical in nature. Both theorems say:
given a map from a product of structured objects to a codomain,
imposing compatibility with each factor (faithfulness/Pareto),
locality (stability/IIA), and non-degeneracy
(decisiveness/non-dictatorship) yields a contradiction. The
difference is in what prevents joint satisfaction:
\emph{underspecification} (multiple models fit the same data) versus
\emph{preference diversity} (multiple voters hold conflicting
orderings).

%% ===================================================================
\subsection{Implications}
\label{sec:arrow-implications}

\paragraph{Independence, not generalization.}
The explanation impossibility is not a generalization of Arrow's
theorem. The two results are logically independent: neither implies
the other. Their shared abstract structure reflects a common
pattern---the impossibility of aggregation under conflicting
locality and non-degeneracy constraints---not a subsumption
relationship. Any claim that one ``contains'' the other would require
either fabricating a Rashomon property where IIA prevents one from
existing, or fabricating a non-dictatorship axiom where the
explanation framework provides none.

\paragraph{The query-relative bridge.}
The identification IIA = query-relative stability could not have been
stated without the query parameterization of the explanation
impossibility (Section~\ref{sec:arrow-setup};
Lean: \texttt{QueryRelative.lean}). The standard
(non-query-relative) framework has a single global incompatibility
relation, which does not capture IIA's pairwise structure. The
query-relative generalization is therefore not merely a technical
refinement: it is the precise bridge between the two frameworks.

\paragraph{Toward a unified categorical framework.}
A unified framework subsuming both theorems would need to abstract
over the \emph{source of impossibility}: Rashomon (many-to-one
observation map) versus preference diversity (many voters with
conflicting orderings). In both cases, a structured input space
admits conflicting ``local views'' that no global map can reconcile.
Formalizing this common pattern---perhaps as a sheaf-theoretic
obstruction or a failure of a limit in a suitable category---is a
direction for future work.

\paragraph{Relation to Okasha (2011).}
\citet{okasha2011theory} applied Arrow's theorem \emph{to} theory
choice, arguing that Kuhnian criteria (accuracy, simplicity,
scope, fruitfulness, consistency) play the role of voters, and
scientific theories play the role of alternatives. This is a
different exercise from ours. Okasha asks: does Arrow's theorem
constrain theory choice? We ask: how does Arrow's theorem
\emph{compare structurally} with the explanation impossibility?
Okasha's analysis uses Arrow as a tool; ours uses Arrow as a
point of comparison. The findings are complementary: Okasha shows
that preference aggregation is relevant to philosophy of science;
we show that the explanation impossibility, while structurally
parallel, operates in a complementary regime where IIA is absent
and Rashomon is present.
